% Part: first-order-logic
% Chapter: model-theory
% Section: theory-of-m

\documentclass[../../../include/open-logic-section]{subfiles}

\begin{document}

\olfileid[fr]{mod}{bas}{thm}

\section{La théorie d'!!a{structure}}

Toute !!{structure}~$\Struct{M}$ rend certains !!{sentence}s vrais
et d'autres faux. L'ensemble des !!{sentence}s qu'elle rend vrais
s'appelle sa \emph{théorie}. Cet ensemble est effectivement une
théorie : tout énoncé qui en découle doit être vrai dans chacun de
ses modèles, notamment dans~$\Struct{M}$.

\begin{defn}
  Étant donné !!a{structure}~$\Struct M$, la \emph{théorie}
  de $\Struct{M}$ est
  l'ensemble $\Theory{M}$ des !!{sentence}s vrais dans $\Struct{M}$,
  c'est-à-dire $\Theory{M} = \Setabs{!A}{\Sat{M}{!A}}$.
\end{defn}

Nous employons aussi le terme « théorie » de façon informelle pour
désigner des ensembles d'!!{sentence}s dotés d'une interprétation
visée, qu'ils soient ou non clos par conséquence déductive.

\begin{prop}
Pour toute $\Struct{M}$, $\Theory{M}$ est complète.
\end{prop}

\begin{proof}
Pour tout !!{sentence}~$!A$, on a soit $\Sat{M}{!A}$, soit
$\Sat{M}{\lnot !A}$; ainsi, soit $!A \in \Theory{M}$,
soit $\lnot !A \in \Theory{M}$.
\end{proof}

\begin{prop}\ollabel{prop:equiv}
  Si $\Struct{N} \models !A$ pour tout $!A \in \Theory{M}$,
  alors $\Struct{M} \elemequiv \Struct{N}$.
\end{prop}

\begin{proof}
Puisque $\Sat{N}{!A}$ pour tout $!A \in \Theory{M}$, on a
$\Theory{M} \subseteq \Theory{N}$. Si $\Sat{N}{!A}$, alors
$\Sat/{N}{\lnot !A}$; donc $\lnot !A \notin \Theory{M}$.
Comme $\Theory{M}$ est complète, $!A \in \Theory{M}$.
Ainsi, $\Theory{N} \subseteq \Theory{M}$, d'où
$\Struct{M} \elemequiv \Struct{N}$.
\end{proof}

\begin{rem}\ollabel{remark:R}
  Considérons $\Struct{R} = \langle\Real, <\rangle$, la !!{structure}
  dont le !!{domain} est l'ensemble $\Real$ des nombres réels, dans
  le !!{language} ne comportant qu'un !!{predicate} binaire,
  interprété par la relation $<$ sur les réels. La !!{structure}
  $\Struct{R}$ est !!{nonenumerable}. Cependant, puisque
  $\Theory{R}$ est manifestement cohérente, le théorème de
  L\"owenheim--Skolem lui donne un modèle !!{enumerable},
  disons $\Struct{S}$; d'après la \olref{prop:equiv}, on a donc
  $\Struct{R} \equiv \Struct{S}$. Puisque $\Struct{R}$ et
  $\Struct{S}$ ne sont pas isomorphes, la réciproque du
  \olref[iso]{thm:isom} est fausse en général.
\end{rem}

\end{document}
